\section{Logics and prior complexity bounds concerning them}

We fix a vocabulary $\vocab$ which consists of $n$ unary predicates $U_1, \ldots, U_n$ and $m$ binary predicates $R_1, \ldots, R_m$. All of our logics will consist of formulas with at most two variables, $x$ and $y$.
We start by reviewing \emph{two variable logic with counting} ($\Ctwo$), built up from atomic formulas $U_i(x)$, $ U_i(y)$, $R_i(x, y)$, and $R_i(y, x)$
via the Boolean operators and the quantifier constructs
$\rho(x) := \exists^{\geq k} y\ \phi(x,y)$,
where $\phi$ is a $\Ctwo$ formula and $k$ is a natural number, and similarly with the role of $x$ and $y$ swapped.

Our semantics will always be over \emph{finite structures} interpreting the vocabulary $\vocab$.
The formula $\rho(y)$ above holds in a model $M$ at an element $u$ for $x$ when there are least $k$ elements $v$ such that $M \models \phi(u, v)$.
The case $k = 1$ represents standard existential quantification.

\emph{Guarded two variable logic with counting} ($\GCtwo$), restricts $\Ctwo$ by requiring quantification to be guarded by an atom. The quantifier rule now includes the forms:
\begin{equation*}
    \exists^{\geq k} y\ R_i(x, y) \land \phi(x,y),
    \quad
    \exists^{\geq k} y\ R_i(y,x) \land \phi(x,y),
    \quad \text{or} \quad
    \exists^{\geq k} x\ ~ U_i(x) \land \phi(x).
\end{equation*}


The logic $\Ptwo$ replaces $\Ctwo$'s counting quantifier with a \emph{Presburger quantifier}:
\begin{equation*}
    \rho(x)\ :=\ \left(
    \sum_{t \in \intsinterval{m}} \kappa_{t} \cdot \presby{ \phi_t(x,y)}  \right)
    \ \circledast\ \delta,
\end{equation*}
where $\kappa_t$ are integers,
$\circledast$ is an inequality or equality, and $\delta$ is an integer.
In $\Ptwo$ we will also allow \emph{unguarded quantification over unary formulas}: if $\phi(x)$ is a formula, so is
$\exists x\ \phi(x)$.
For example,
$\exists x \left(\presby{R(x,y)} - 3 \cdot \presby{R(y,x)} = 0\right)$
is a sentence of $\Ptwo$ that holds in a directed graph  when there is an element that has 3 times as many outgoing edges as incoming edges.

The logic $\GPtwo$ (\emph{guarded two variable logic with Presburger quantifiers}) restricts $\Ptwo$ by requiring each $\phi_i$ in a Presburger quantifier to be guarded, although still allowing unguarded unary quantification.
Note that global cardinality constraints can be expressed in $\Ptwo$, but not in $\GPtwo$.
Following prior papers \cite{twovarpres}, we do not allow a quantifier that tests the modulus of a linear combination of
cardinalities of one dimensional sets, and thus our quantifiers are in a sense weaker than Presburger arithmetic.
A \emph{Presburger modulus constraint} is as above, but allowing $\circledast\ \delta$ to be replaced by 
$(= m \mod n)$, for natural numbers $m, n$ with $m < n$.
Likewise, $\Ctwoplusglobal$ can be extended to allow modulus constraints. 
In the appendix we show that \emph{all of of our results on $\GPtwo$ and $\Ctwoplusglobal$ also hold for the extension with modulus constraints}; in the body of the paper we focus on plain $\GPtwo$ and $\Ctwoplusglobal$ in the statements. 


Given a unary formula $\phi(x)$ in some logic, and a finite model $M$, we let $|\phi|_M$ be the cardinality of $\phi$ in $M$; we extend this notation to a tuple $\langle \phi_1(x) \ldots \phi_k(x) \rangle$.
Given such a tuple, its \emph{spectrum} is the set of vectors of numbers of the form  $| \phi_1(x) \ldots \phi_k(x) |_M$. We say that such a spectrum is \emph{semi-linear} if it can be represented as a set of linear inequalities. When we talk about \emph{computing the spectrum} for a formula, we mean computing a semi-linear representation. 
If we have such a representation, we can clearly check whether an individual $\phi_i(x)$ is satisfiable, simply by seeing if the corresponding set of constraints is feasible. Thus, in cases, where the spectrum is known always to have such a representation, computing the spectrum can be seen as a generalization of the satisfiability pboelm.


The positive results in this paper, as well as the prior bounds, are summarized in Table \ref{tab:results}.
In the table, Sat Logic refers to the satisfiability problem for formulas of the logic over finite structures, while
Spectrum refers to a bound on the size of a representation of the spectrum, again over finite structure. An entry of the form $X$-c indicates that the problem is  complete for complexity class $X$.
\begin{tiny}
\begin{table}[!ht] \
\centering
\resizebox{\columnwidth}{!}{%
\begin{tabular}{l|l|l}
\hline
& Sat Logic     &   Spectrum Size Bound \\
\hline
$\GCtwo$ & $\exptime$-c \cite{gctwocomplexity} & {\bf $\expp$} (formerly
  $\twoexpp$) \\ %\cite{ctwocomplexity}
$\Ctwo$ & $\nexptime$-c \cite{ctwocomplexity}  & {\bf $\expp$, Thm \ref{thm:spectrasemilinear}} (formerly  
$\twoexpp$ \cite{ultimatelyperiodic}) \\
$\GPtwo$ & {\bf $\exptime$-c, Thm \ref{thm:gptwoexp}} (formerly $\exptime$-hard, in $
\threenexp$ \cite{percentage}) & Open \\
$\Ptwo$ &  Undecidable \cite{percentage} & No size bound \\
\hline
\end{tabular}
}%resizebox
\caption{Complexity results, with those new to this work in bold.}
\label{tab:results}
\end{table}
\end{tiny}


We give results on the spectrum, but they have applications to decidability of  richer logics. $\Ctwo$ can be extended with \emph{global unary cardinality constraints}: these are of
the form
$\rho := 
\left(\sum_{t \in \intsinterval{n}}
\kappa_t \cdot |\phi_t(x)|\right)\ \circledast\ \delta$,
where $\phi_t$ are $\Ctwo$ formulas in one free variable,
$\kappa_t$ are integers,
$\circledast$ is an inequality or equality,
and $\delta$ is an integer.
$\rho$ is a sentence, and we define the semantics only for finite structures $M$.
Such $M$ satisfies $\rho$ if,
letting $s_t$ be the number of elements in
$M$ satisfying $\phi_t(x)$, we have 
$\left(\sum_{t \in \intsinterval{n}} \kappa_t \cdot s_t\right)\ \circledast\ \delta$.
We write $\Ctwoplusglobal$ for the extension of $\Ctwo$ with global cardinality constraints.
As we explain later, computation of the spectrum provides an approach to get tight bounds on these extended logics. In the case where the global constraints are in Presburger arithmetic, these can be seen as an alternative to the technique of \cite{rudolphcardinality}, which obtains tight bounds for the logic through reduction to satisfiability of $\Ctwo$ itself.